\section{Priced sectors and an exterior cutoff}\label{sec:pricing}

The price per electron permits insertion and deletion comparisons without a
constraint on the exterior particle number.  We first collect the screening
estimates that remain applicable when the energy above the unpriced minimum
grows.  Throughout, write
\[
 K(x)=\frac1{|x|},\qquad V=ZK,\qquad \nu=Z\delta_0,\qquad
 D(f)=\frac12\iint_{\R^3\times\R^3}
                  \frac{f(x)f(y)}{|x-y|}\,dx\,dy,
\]
and use the constants
\[
 T=\frac3{10}(3\pi^2)^{2/3},\qquad
 k=\left(\frac{5T}{3}\right)^{-3/2}.
\]
These are the Thomas--Fermi constants for two spin states and kinetic energy
$-\Delta/2$.  For a fixed integer $Z$, abbreviate $E_n=E_Z(n)$, put $E_0=0$,
and define
\[
 E=\inf_{n\ge0}E_n,\qquad
 P(\lambda)=\min_{n\ge0}(E_n+\lambda n)\quad(\lambda>0).
\]
The existence of this minimum follows below.  The asymptotic regime used until
\cref{sec:fixed-sectors} is
\begin{equation}\label{eq:scaling-regime}
 \lambda=s^{-4},\qquad s\longrightarrow0,
 \qquad Zs^3\longrightarrow\infty.
\end{equation}
Every assertion concerning a minimizing sector applies to every choice of
minimizing index at every system in this regime.

If $\psi$ is a normalized finite-sector form-domain vector, we call it an
\emph{unshifted $\delta$-state} when its form energy is at most $E+\delta$,
where $\delta\ge0$.  Its position law is $|\psi|^2$, with the spin variables
summed unless they are recorded; its one-particle density is denoted by
$\rho_\psi$.  Expectations may also include the auxiliary random labels
specified below.  For a real number $v$, write $v_+=\max(v,0)$ and
$v_-=\max(-v,0)$.  Constants $C,c>0$ can change from line to line.  Unless a
dependence is displayed, they depend only on this fixed model and on the fixed
cutoff and kernel conventions.

\subsection{Imported estimates and their hypotheses}

\begin{lemma}[Sector and energy estimates]\label{lem:sector-inputs}
The sequence $E_n$ is nonincreasing.  If $E_n<E_{n-1}$, there is a normalized
ground state in sector $n$, and $n\le2Z+1\le3Z$.  There is a largest such
index $N_*$, with $Z\le N_*\le3Z$, and
\[
 E=E_{N_*}=E_n\quad(n\ge N_*).
\]
There are universal $C,Z_0$ such that
\begin{equation}\label{eq:neutral-offset}
 E_Z(Z)\le E+C\qquad(Z\ge Z_0).
\end{equation}
If $\Psi$ is a normalized sector ground state and $F$ is a bounded real
symmetric Lipschitz function of the spatial configuration, then
\begin{equation}\label{eq:ground-multiplier}
 \langle F\Psi,H_{Z,n}F\Psi\rangle
       -E_n\|F\Psi\|_2^2
       =\frac12\E_\Psi|\nabla F|^2,
\end{equation}
where the left side is interpreted as a quadratic form and the gradient
includes all spatial coordinates.  Finally, every unshifted $\delta$-state
in a sector $n\le3Z$ satisfies
\begin{equation}\label{eq:global-density-input}
 \int_{\R^3}\rho_\psi^{5/3}\le C(Z^{7/3}+\delta).
\end{equation}
\end{lemma}

\begin{proof}
The sector assertions and the multiplier identity are
\cite[\FSectors]{Foundation}; the density estimate is
\cite[\FGlobalDensity]{Foundation}.  The stronger estimate
$E_Z(Z-1)\le E+C$ is \cite[\FNeutralOffset]{Foundation}.
Monotonicity gives \cref{eq:neutral-offset}.  All these results use exactly
the Hamiltonian and the two-spin convention fixed here.
\end{proof}

A one-particle square cut consists of nonnegative bounded-gradient Lipschitz
functions $o,c$ with $o^2+c^2=1$.  Conditional on the positions, assign the
particles independently to the out and core groups with probabilities $o^2$
and $c^2$.  The \emph{cut data} record the out labels, positions, and spins.
For a finite sequence of cuts, stage labels are recorded as well.  Conditional
core states are obtained by the corresponding products of cut factors, so
antisymmetry is preserved within each group.  This is the localization
convention of \cite[\FLocalization]{Foundation}.

The local estimate needed for the later patch comparison has a quantitative
condition on these cuts.  Fix, once and for all,
\[
 L=10^5,\qquad A=40,\qquad a_y=|y|/L,\qquad B_y=B(y,a_y)
       \quad(y\ne0).
\]
For an unshifted $\delta$-state put
\[
 \begin{gathered}
 m_y=\max(a_y^{-3},1)+\sqrt{\delta a_y},\qquad
 J_y=V(y)-\sum_{|x_i-y|\ge Aa_y}\frac1{|y-x_i|},\\
 P_* =\max\left(1,\sup_{y\ne0}
                 \frac{\E\bigl[\#\{i:x_i\in B_y\}^2\bigr]}{m_y^2}\right).
 \end{gathered}
\]
The number $P_*$ is finite for each finite system, since the count is at most
the sector particle number and $m_y\ge1$.

\begin{lemma}[Local fields for arbitrary offsets]\label{lem:local-field-input}
For every normalized unshifted $\delta$-state of finite particle number,
$P_*\le C$.  Fix $y\ne0$ and $C_0\ge1$.  The data $X$ may be trivial.
Otherwise, suppose that the initial cut, or finite sequence of cuts
$(o_h,c_h)$, has total out support covered by at most $C_0$ cells $B_z$ with
\[
 C_0^{-1}\le a_z/a_y\le C_0,
 \qquad
 \sum_h\bigl(|\nabla o_h|^2+|\nabla c_h|^2\bigr)
       \le\sum_z D_z^2a_z^{-2}\ind_{B_z},
 \qquad
 \frac{D_z^2}{a_zm_z}\le C_0\sqrt{P_*}.
\]
Then, for the complete cut data,
\begin{equation}\label{eq:local-field-input}
 \left\|(\E[J_y\mid X])_+\right\|_2
       \le C(C_0)\frac{m_y}{a_y}.
\end{equation}
The constants are independent of $Z$, the sector, the state, $y$, and
$\delta$, apart from the displayed dependence through $m_y$.
\end{lemma}

\begin{proof}
This is \cite[\FLocalField]{Foundation}.  Its offset is measured from the
minimum $E$ over all sectors, just as here; it has no restriction that the
offset stay bounded along a sequence.  In all applications below we verify
the stronger cut condition $D_z^2/(a_zm_z)\le C_0$.
\end{proof}

\begin{lemma}[Annular screening estimates]\label{lem:screening-inputs}
Let $\psi$ be a normalized unshifted $\delta$-state of finite particle number.
For fixed $0<\alpha<\beta<\infty$ and every $u>0$,
\begin{equation}\label{eq:annular-count}
 \left\|\#\{i:\alpha u\le|x_i|\le\beta u\}\right\|_2
   \le C_{\alpha,\beta}
           \bigl(\max(u^{-3},1)+\sqrt{\delta u}\bigr).
\end{equation}
Suppose in addition that $n\le3Z$, and set
\[
 B_{u/8}(x)=V(x)-\sum_{|x_i|<u/8}\frac1{|x-x_i|}.
\]
Let $Y$ be trivial or be the data of a square cut whose out support lies in
$\{u/4\le|x|\le8u\}$ and whose squared gradient sum is at most $C/u^2$.
There are conditional versions harmonic on $u/4<|x|<6u$ for which
\begin{equation}\label{eq:conditional-inner-field}
 \left\|\sup_{u/2\le|x|\le4u}
              \bigl(\E[B_{u/8}(x)\mid Y]\bigr)_+\right\|_2
 \le C\left(\frac{\max(u^{-3},1)}u+\sqrt{\frac\delta u}\right).
\end{equation}
The constant in the second estimate may depend on the fixed gradient bound.
Both estimates allow every finite $\delta$.
\end{lemma}

\begin{proof}
The count estimate is \cite[\FAnnularCount]{Foundation}.  Boundary spheres
have zero probability because the one-particle density is absolutely
continuous, so open or closed annuli give the same estimate.  For the field
bound, use \cite[\FInnerField]{Foundation} with inner cutoff $h=u/8$, broad
annulus $u/4<|x|<6u$, and smaller annulus $u/2\le|x|\le4u$.
The series in that lemma is
\[
 \frac{\max(u^{-3},1)+\sqrt{\delta u}}u
 +\sum_{j\ge0}
       \frac{\max((2^jh)^{-3},1)+\sqrt{\delta\,2^jh}}
            {\max(u,2^jh)}.
\]
With $h=u/8$, geometric summation bounds it by the right side of
\cref{eq:conditional-inner-field}, also when $u<1$.  The permitted cut
support and gradient bound satisfy the fixed-annulus hypotheses of that
lemma.  The conditional harmonic versions follow by integration against
the conditional law: for each fixed finite system the inner sources are
uniformly separated from the annulus, so differentiation under the integral
is valid.  No minimizing property is required in either estimate.
\end{proof}

\subsection{The rough charge deficit}

\begin{lemma}\label{lem:priced-rough}
The minimum $P(\lambda)$ is attained.  Along \cref{eq:scaling-regime}, every
minimizing index $n$ is eventually nonzero and strictly bound.  For a
normalized ground state $\Psi$ in that sector, put $D_0=E_n-E$.  Uniformly
over all these choices,
\begin{equation}\label{eq:rough-deficit}
 |Z-n|\le Cs^{-3},\qquad 0\le D_0\le Cs^{-7}.
\end{equation}
\end{lemma}

\begin{proof}
By \cref{lem:sector-inputs}, $E_n$ is constant for $n\ge N_*$.  Since
$\lambda>0$, the priced objective is strictly increasing thereafter, so
its minimum is attained.  If a minimizing index is positive, comparison
with $n-1$ gives
\[
 E_{n-1}-E_n\ge\lambda>0.
\]
It is therefore strictly bound, has a ground state, and satisfies $n\le3Z$.
A normalized exponential one-electron trial of radius $Z^{-1}$ has energy
$-Z^2/2$.  Moreover,
\[
 \frac\lambda{Z^2}=\frac{s^2}{(Zs^3)^2}\longrightarrow0.
\]
Thus $P(\lambda)<0$ eventually, excluding the vacuum.

The neutral sector is an admissible competitor.  By
\cref{eq:neutral-offset},
\begin{equation}\label{eq:priced-offset}
 0\le D_0\le C+s^{-4}(Z-n).
\end{equation}
For a fixed inner configuration, the spherical average of $B_{s/8}$ on
$|x|=s$ equals $(Z-\#\{|x_i|<s/8\})/s$.  Its expectation is at least
$(Z-n)/s$.  Apply \cref{eq:conditional-inner-field} with trivial data and
offset $D_0$ to obtain, for small $s$,
\[
 Z-n\le C\bigl(s^{-3}+\sqrt{D_0s}\bigr).
\]
If $d=Z-n\ge0$ and $z=s^3d$, then \cref{eq:priced-offset} gives
$D_0s^7\le Cs^7+z$, whence
\[
 z\le C\bigl(1+\sqrt{z+Cs^7}\bigr).
\]
This bounds $z$.  If $d<0$, \cref{eq:priced-offset} instead gives
$d\ge-Cs^4$.  These two cases prove the deficit estimate, and substitution
in \cref{eq:priced-offset} proves the bound for $D_0$.
\end{proof}

\subsection{A cutoff at one fixed exterior scale}

\begin{lemma}\label{lem:exterior-cutoff}
There is a sufficiently large fixed $R_0$ such that, for every sequence of
minimizing sectors and corresponding ground states in
\cref{eq:scaling-regime},
\begin{equation}\label{eq:exterior-cutoff}
 s^3\E_\Psi\#\{i:|x_i|>R_0s\}\longrightarrow0.
\end{equation}
\end{lemma}

\begin{proof}
Choose a smooth $0\le\chi\le1$ supported in $[1/2,4]$ and equal to one on
$[1,2]$.  At shell scale $u>0$, write
\[
 w_i=\chi(|x_i|/u)^2,\quad S_u=\sum_iw_i,\quad m_u=\E S_u,
 \quad N_u=\E\#\{i:u/2\le|x_i|\le4u\}.
\]
Test the weak eigen-equation with $S_u\Psi$.  In its $i$th summand the other
coordinates have energy at least $E_{n-1}$.  Integration by parts gives for
the selected-coordinate kinetic term
\[
 \frac12\operatorname{Re}\int\nabla_i\overline\Psi
                         \cdot\nabla_i(w_i\Psi)
 =\frac12\int w_i|\nabla_i\Psi|^2
       -\frac14\E\Delta_iw_i
 \ge-\frac14\E\Delta_iw_i.
\]
Retain only the repulsions from particles in $|x_l|<u/8$; the selected
particle lies outside that ball whenever $w_i>0$.  Since
$E_{n-1}-E_n\ge s^{-4}$, we obtain
\begin{equation}\label{eq:shell-eigen-equation}
 s^{-4}m_u\le Cu^{-2}N_u
                 +\E\sum_iw_iB_{u/8}(x_i).
\end{equation}
All these tests are bounded symmetric form-domain multipliers for each
fixed system.

Suppose $m_u>0$.  Bias the position law by $S_u/m_u$, equivalently use the
normalized state
\[
 \Psi_u=\sqrt{S_u/m_u}\,\Psi.
\]
The function $\sqrt{S_u}$ is the Euclidean norm of the vector of cutoffs.
It is Lipschitz, including at its zeros, and
\[
 |\nabla\sqrt{S_u}|^2
       \le Cu^{-2}\#\{i:u/2\le|x_i|\le4u\}.
\]
The multiplier identity \cref{eq:ground-multiplier} therefore makes
$\Psi_u$ an unshifted state of offset at most
\begin{equation}\label{eq:shell-bias-offset}
 \delta_u=D_0+Cu^{-2}N_u/m_u.
\end{equation}

In both laws take the same auxiliary square cut, full out on
$u/2\le|x|\le4u$, supported in $u/4\le|x|\le8u$, with gradient bounded by
$C/u$.  Its data $Y$ determine $S_u$ and every position with positive
weight.  Biasing by the $Y$-measurable variable $S_u/m_u$ does not change
the conditional raw law on $S_u>0$.  Consequently the two conditional
harmonic fields there are the same.  Let
\[
 M(Y)=\sup_{u/2\le|x|\le4u}
            \bigl(\E[B_{u/8}(x)\mid Y]\bigr)_+
\]
be this common conditional supremum.  Conditioning first on $Y$ gives
\[
 \E\sum_iw_iB_{u/8}(x_i)
       \le\E[S_uM(Y)]
       =m_u\E_{\Psi_u}M(Y).
\]
The equality of conditional fields is used only on positive-bias data;
zero-weight summands are omitted.  Their harmonic continuity justifies
evaluation at the recorded positions.  Applying
\cref{eq:conditional-inner-field} to $\Psi_u$, with
\cref{eq:shell-bias-offset}, now yields
\begin{align}
 \E\sum_iw_iB_{u/8}(x_i)
 &\le Cm_u\left(\frac{\max(u^{-3},1)}u
                            +\sqrt{D_0/u}\right)
       +Cu^{-3/2}\sqrt{m_uN_u}.
       \label{eq:shell-field-bound}
\end{align}
If $m_u=0$, the resulting shell estimate is immediate as well.

By \cref{eq:rough-deficit}, for $u\ge R_0s$,
\[
 s^4\frac{\max(u^{-3},1)}u\le R_0^{-4}+s^4,
 \qquad
 s^4\sqrt{D_0/u}\le C\sqrt{s/u}\le CR_0^{-1/2}.
\]
Choose $R_0$ large and then $s$ small.  The terms multiplying $m_u$ in
\cref{eq:shell-field-bound} can be absorbed in
\cref{eq:shell-eigen-equation}.  As $m_u\le N_u$, the mixed term is at
most $Cu^{-3/2}N_u$, and hence
\[
 m_u\le Cs^4(u^{-2}+u^{-3/2})N_u.
\]
The annular count bound and \cref{eq:rough-deficit} also give
\[
 N_u\le C\bigl(u^{-3}+1+s^{-7/2}\sqrt u\bigr).
\]
The shells on which the weights equal one, with
$u=2^jR_0s$, $j\ge0$, cover the required exterior.  Tonelli's theorem and
the two preceding inequalities give
\begin{align*}
 s^3\E\#\{|x_i|>R_0s\}
 &\le C s^7\sum_{j\ge0}
       (u^{-2}+u^{-3/2})(u^{-3}+1+s^{-7/2}u^{1/2})
       \bigg|_{u=2^jR_0s}\\
 &\le C_{R_0}
       \bigl(s^2+s^{5/2}+s^5+s^{11/2}\bigr)
       \longrightarrow0.
\end{align*}
Every power of $u$ in this sum is negative.  Thus the estimate includes
arbitrarily distant shells, with the same fixed $R_0$.
\end{proof}
